\documentclass[11pt,a4paper]{article}
\usepackage[margin=24mm]{geometry}
\usepackage[T1]{fontenc}
\usepackage{lmodern}
\usepackage{amsmath,amssymb,amsthm,mathtools}
\usepackage[hidelinks]{hyperref}
\usepackage{bookmark}
\setlength{\emergencystretch}{2em}
\newtheorem{lemma}{Lemma}
\newtheorem{proposition}{Proposition}
\newcommand{\Q}{\mathbb Q}
\newcommand{\R}{\mathbb R}
\newcommand{\C}{\mathbb C}
\DeclareMathOperator{\Cl}{Cl}
\DeclareMathOperator{\Supp}{Supp}
\DeclareMathOperator{\coeff}{coeff}
\title{Local $\Q$-Triviality of the Canonical Divisor}
\author{}
\date{}
\hypersetup{pdftitle={Local Q-Triviality of the Canonical Divisor},
  pdfsubject={Statements following the handwritten outline; proofs to be added}}

\begin{document}
\maketitle
\vspace{-2.2em}

Let $X$ be a normal projective variety over $\C$, and let $f:X\to X$
be an int-amplified endomorphism. Set
\[
 P:=\frac{f_*}{\deg f}.
\]
For a finite surjective morphism $h$, write $R_h$ for its ramification
divisor. For birational morphisms, we use the Weil pullback.
For a prime divisor $F$ on a normal proper birational model $p:Y\to X$,
write $c_X(F)=p(F)$ and
\[
 a^+(F,X,0):=\coeff_F\bigl(K_Y+p^*(-K_X)\bigr).
\]
Canonical divisors are chosen compatibly.

\begin{lemma}\label{lem:polynomial}
There exists a polynomial $q(t)\in\Q[t]$ such that
\[
 q(P)=0\quad\text{on }\Cl(X)_{\Q},
 \qquad q(1)\ne0.
\]
\end{lemma}
% Proof of Lemma 1: to be added by the user.
% \begin{proof}
% \end{proof}

\begin{lemma}\label{lem:decomposition}
There exist an integer $N\ge1$ and rational numbers
$a_1,\ldots,a_N,b_1,\ldots,b_N$ such that
\[
 K_X\sim_{\Q}\sum_{i=1}^N a_iP^iR_f
     =\sum_{i=1}^N b_iP^iR_{f^i}.
\]
\end{lemma}
% Proof of Lemma 2: to be added by the user.
% \begin{proof}
% \end{proof}

\begin{lemma}\label{lem:finite-base-change}
Consider a commutative diagram of normal varieties
\[
 \begin{array}{ccc}
  Y'&\xrightarrow{\ f'\ }&Y\\[2pt]
  {\scriptstyle p'}\downarrow&&\downarrow{\scriptstyle p}\\[2pt]
  X_s&\xrightarrow{\ f\ }&X_t,
 \end{array}
\]
where $f,f'$ are finite surjective and $p,p'$ are proper birational.
For every Weil $\R$-divisor $D$ on $X_t$, one has
\[
 (f')^*(p^*D)=(p')^*(f^*D).
\]
\end{lemma}
% Source for the statement: Boucksom--de Fernex--Favre,
% The volume of an isolated singularity, Proposition 2.19 and Remark 2.4.
% https://arxiv.org/pdf/1011.2847
% Proof of Lemma 3: to be added by the user.
% \begin{proof}
% \end{proof}

\begin{proposition}\label{prop:avoidance}
Let $F$ be a prime divisor over $X$ with $a^+(F,X,0)=-1$, and set
$Z=c_X(F)$. Then, for every integer $i\ge1$,
\[
 \eta_Z\notin\Supp\bigl((f^i)_*R_{f^i}\bigr).
\]
\end{proposition}
% Proof of Proposition 1: to be added by the user.
% The proof may use the previously established global lower bound
% a^-(E,X,0) >= -1 under the int-amplified hypothesis.
% \begin{proof}
% \end{proof}

\begin{proposition}\label{prop:local-q-triviality}
Let $F$ be a prime divisor over $X$ with $a^+(F,X,0)=-1$, and set
$Z=c_X(F)$. There exists an open neighborhood $U\subseteq X$ of
$\eta_Z$ such that
\[
 K_X|_U\sim_{\Q}0.
\]
\end{proposition}
% Proof of Proposition 2: to be added by the user.
% \begin{proof}
% \end{proof}

\end{document}
